\section{The scalar and path-valued pointwise errors coincide}
\label{sec:v73-scalar-path}
The full-source estimate should be used for all bounded-Lipschitz
path tests on a common version, rather than proved anew for each
numerator. The inherited positive path-remainder theorem gives an
exact comparison of the two asymptotic errors. This comparison keeps
the scalar endpoint as the target; it is not an assertion that the
error has already vanished.

Let $\mathcal C=C([0,1],\R^4)$ and use the bounded-Lipschitz class
$\mathcal F$ of Section~\ref{sec:v52-path-inversion}, including the
constant function one. At a fixed roof put
\[
 \mathbf P(t)=m^2p(t)\mathsf Q_t,\qquad P(t)=\mathbf P(t)(1),
 \qquad G(t)=\mathcal L_{m,R}(k_1,k_2,t,n).
\]
The reference path law is $\mathsf W_R$, with the same covariance
and finite arithmetic coefficient as in
Theorem~\ref{thm:v52-path-remainder}. For each fixed $B$ let
$\mathbf b_B$ be the kernel obtained from the original physical
weight $1-H^{\varepsilon(B)}$. Its mass is $m^2b^{\varepsilon(B)}$.
All kernels are those of Lemma~\ref{lem:v73-common-kernel}.

Define the original central roof set and two ordered errors by
\begin{align}
 \mathcal C_{m,R,n,k}(M)&=\{t:|Z_{m,R}(k,t,n)|\le M\},\notag\\
 E_{\mathrm s}(M)&=\limsup_m\sup_{R,n,k}
       \operatorname*{ess\,sup}_{t\in\mathcal C_{m,R,n,k}(M)}
                                |P(t)-G(t)|,\notag\\
 E_{\mathrm p}(M)&=\limsup_m\sup_{R,n,k}
       \operatorname*{ess\,sup}_{t\in\mathcal C_{m,R,n,k}(M)}
                       \|\mathbf P(t)-G(t)\mathsf W_R\|_{\mathrm{BL}^*}.
                                                   \label{eq:v73-errors}
\end{align}
An empty target contributes zero. In particular $E_{\mathrm s}(M)$
is the original $\mathcal E_M$, not an interval-averaged error.

\begin{theorem}[Equality of the ordered pointwise errors]
\label{thm:v73-scalar-path-equality}
On the hypotheses and conclusions of the inherited
Theorem~\ref{thm:v52-path-remainder}, for every $M<\infty$,
\begin{equation}\label{eq:v73-scalar-path-equality}
                    E_{\mathrm p}(M)=E_{\mathrm s}(M).
\end{equation}
Consequently vanishing of the original scalar pointwise error is
equivalent to vanishing of the full path-numerator error in the
bounded-Lipschitz dual norm at the same roofs and exact labels.
This statement does not assert total-variation convergence on path
space.
\end{theorem}
\begin{proof}
For fixed $B$ write
$\mathbf E_B=\mathbf P-G\mathsf W_R-\mathbf b_B$ and let
$e_B=\|\mathbf E_B\|_{\mathrm{BL}^*}$. Positivity gives
$\|\mathbf b_B\|_{\mathrm{BL}^*}=\mathbf b_B(1)$. The constant
test one and the triangle inequality therefore give pointwise,
on the common versions,
\begin{equation}\label{eq:v73-pointwise-norm-comparison}
 |P-G|\le\|\mathbf P-G\mathsf W_R\|_{\mathrm{BL}^*}
                                      \le |P-G|+2e_B.
\end{equation}
Indeed the middle term is at most $\mathbf b_B(1)+e_B$, whereas
$|P-G|=|\mathbf b_B(1)+\mathbf E_B(1)|
\ge\mathbf b_B(1)-e_B$. The inherited theorem bounds the collision
limsup of the global essential supremum of $e_B$ by
$CB^{-1/192}$. Restricting to the original central set, taking the
collision limit at fixed $B$, and then sending $B$ to infinity
proves \eqref{eq:v73-scalar-path-equality}. If the scalar error is
infinite the lower inequality already gives equality in the extended
sense. No interchange of a supremum over uncountably many path
versions is made: one common kernel and a countable norming class
were fixed first.
\end{proof}

The input to this theorem includes the inherited cylinder inversion,
variation tightness and endpoint budget in modules 110 and 112.
The equality neither reconstructs that spectral chain nor replaces
its specialist audit. It shows that, once that chain is used as
written, a second independent vanishing-height estimate for the
path numerator is unnecessary.

\begin{corollary}[Central excess and positive arithmetic classes]
\label{cor:v73-central-excess}
For fixed $B,\delta,K$, define
\[
 U_M(B,\delta,K)=\limsup_m\sup_{R,n,k}m^2
  \operatorname*{ess\,sup}_{t\in\mathcal C_{m,R,n,k}(M)}
             [b^{\varepsilon(B)}-KA_\delta b^{\varepsilon(B)}]_+(t).
\]
For $K\ge1$ there is $C_M$ independent of $B,\delta,K$ such that
\begin{align}
 U_M(B,\delta,K)&\le E_{\mathrm s}(M)+C_MB^{-1/192},\notag\\
 E_{\mathrm s}(M)&\le U_M(B,\delta,K)
                         +C_M(1+K)B^{-1/192}.\label{eq:v73-central-sandwich}
\end{align}
Thus if $K(B)B^{-1/192}\to0$, the central excess has the same
outer limit as the scalar and path errors, whenever these errors
are finite. In particular their vanishing is equivalent. The paired
flux estimate of Theorem~\ref{thm:v73-paired-budget} is a sufficient
bound for this excess with $K=2$.

On any of these central targets with $G\ge d>0$, vanishing of
$E_{\mathrm s}(M)$ gives, uniformly in essential supremum,
\[
 d_{\mathrm{BL}}(\mathsf Q_t,\mathsf W_R)\longrightarrow0,
 \qquad P/G\longrightarrow1,\qquad G/P\longrightarrow1.
\]
The unnormalized error comparison includes zero arithmetic classes;
none of the displayed ratios or conditional laws is assigned there.
On each original roof window contained in this central target with
the same reference floor, the forward normalized arithmetic likelihood
of Section~\ref{sec:v54-forward-likelihood} also tends to one in
essential supremum.
\end{corollary}
\begin{proof}
The scalar test in the positive path-remainder theorem gives
$P-G=m^2b+\mathbf E_B(1)$. Hence the ordered central height of
$b$ is at most $E_{\mathrm s}(M)+C_MB^{-1/192}$, proving the first
inequality. Conversely
$b\le KA_\delta b+[b-KA_\delta b]_+$, and the fixed-width estimate
$\mathcal H(A_\delta b)\le C\varepsilon(B)^{1/16}$ proves the
second. All widths and reserves are fixed before the collision
limit. This is localization to the target already in the original
pointwise theorem, not deletion of any physical source or label.

For the conditional assertion let
$D=\|\mathbf P-G\mathsf W_R\|_{\mathrm{BL}^*}$ and
$a=|P-G|$. Where $G\ge d$ and $a\le d/2$, $P\ge d/2$ and
\[
 d_{\mathrm{BL}}(\mathsf Q_t,\mathsf W_R)
       \le\frac{D+a}{P}\le\frac{2(D+a)}d,
 \qquad |P/G-1|\le a/d,\qquad |G/P-1|\le2a/d.
\]
Theorem~\ref{thm:v73-scalar-path-equality} gives the required uniform
limits. These are pointwise ratios of the original scalar densities;
no absolute continuity of one path law with respect to another is
asserted.

For the last assertion take a roof window $J$ of positive length
on which $G\ge d$. Write $F_J=\int_J P$, $H_J=\int_J G$,
and put $r=\|P-G\|_{L^\infty(J)}/d$. Then
$|F_J/H_J-1|\le r$. For $r<1$ the original normalized roof
laws satisfy
\[
 \left\|\frac{d\mathbb P}{d\mathbb Q}-1\right\|_\infty
 =\left\|\frac{P/G}{F_J/H_J}-1\right\|_\infty
                         \le\frac{2r}{1-r}.
\]
Thus the forward essential-likelihood consequence requires no
additional numerator estimate after the scalar endpoint. This is
a consequence for the original window-normalized likelihood, not
a replacement of the pointwise theorem by a window law.
\end{proof}

The actual-return path continues to require the inherited collision-to-
return clock transfer on the same positive arithmetic classes. The
new comparison supplies its same-roof collision numerator once the
scalar endpoint is closed; it introduces neither a new clock nor an
unproved conditional normalization. At present the quantitative
paired-flux decay for the unrestricted Lorentz family remains to be
established, so the unconditional scalar and return-path endpoint is
not asserted by this comparison.
\label{end:v73-new-mathematics}
